% !TEX root = ../my-thesis.tex

\lstset{breaklines=true, breakatwhitespace=false}

\chapter{Modellierung und lokale Evaluation}
\label{sec:lokal}

\section{Lokale Untersuchungsumgebung und Systembereitstellung}
\label{sec:lokal:umgebung}

Alle Modellierungen und die lokale Messreihe wurden auf einem MacBook
Pro (2021) mit Apple-M1-Pro-Prozessor (arm64) und 16\,GB
Arbeitsspeicher unter macOS~26.5.1 ausgeführt. Für jedes der
Python-basierten Systeme wurde eine eigene Umgebung angelegt (conda
oder venv), um wechselseitige Abhängigkeitskonflikte auszuschließen. Eine HPC-Ausführung wird damit nicht abgebildet, die lokale Ebene
dient der Modellierung und der Messung unter kontrollierten
Bedingungen. Wie sich die
Systeme auf einem Slurm-basierten HPC-System betreiben lassen,
untersucht Kapitel~\ref{sec:mogon}, auf die Grenzen der lokalen
Betrachtung geht Abschnitt~\ref{sec:diskussion:grenzen} ein.

Die folgenden Absätze beschreiben die Bereitstellung je System in
derselben Reihenfolge: eingesetzte Version und Installationsweg,
zusätzlich benötigte Komponenten und die dafür nötigen
Konfigurationsschritte sowie aufgetretene Fehler.

\paragraph{Nextflow:} Nextflow wurde in Version~25.10.4 über Homebrew
installiert. Als Voraussetzung benötigt Nextflow eine Java-Laufzeit
ab Version~17~\cite{nextflowdocs-install}, die die verwendete
Homebrew-Formel als Abhängigkeit mitinstalliert~\cite{brewnextflow}.
Weitere Komponenten und Konfigurationsschritte waren für den lokalen
Betrieb nicht erforderlich. Bei der Einrichtung traten keine Fehler
auf.

\paragraph{StreamFlow:} StreamFlow wurde in der Vorabversion~0.2.0rc2
eingesetzt. Die Installation erfolgte mit pip direkt aus dem
Quellcode-Repository des Projekts~\cite{streamflowrepo} und wurde auf
den Commit~\texttt{7704a8e} festgelegt. StreamFlow wird als
Python-Paket installiert und setzt Python ab Version~3.10
voraus~\cite{streamflowdocs}. Zusätzliche Komponenten waren nicht
erforderlich. Konfigurationsschritte über die
Installation hinaus und Fehler traten nicht auf.

\paragraph{Merlin:} Merlin wurde in Version~1.13.0 mit pip
installiert. Für die Ausführung benötigt Merlin einen
Nachrichtenvermittler, in dessen Warteschlange Celery die Schritte
als Aufgaben einstellt, die Worker entgegennehmen und
abarbeiten~\cite{merlindocs-celery}. In der lokalen Einrichtung diente Redis als Nachrichtenvermittler und wurde über einen Docker-Container bereitgestellt. Die Verbindungsdaten zum Nachrichtenvermittler stehen in der Datei \texttt{app.yaml} im Nutzerverzeichnis. Redis-Container und Celery-Worker müssen bei der Ausführung aktiv sein, bleiben aber über mehrere Läufe hinweg bestehen und wurden
daher einmal je Sitzung gestartet. Bei der Einrichtung trat ein
Fehler auf: Aktuelle Versionen von \texttt{setuptools} liefern das
von Merlin vorausgesetzte Modul \texttt{pkg\_resources} nicht mehr
mit, sodass das Werkzeug nach der Installation zunächst nicht
startete. Erst nach einem Downgrade von \texttt{setuptools} auf eine
Version kleiner als~81 lief Merlin an.

\paragraph{Pegasus:} Pegasus wurde in Version~5.1.2 installiert und
benötigt HTCondor~\cite{pegasusdocs-install}, das in Version~23.1.0 über
Homebrew bezogen wurde. Die HTCondor-Prozesse wurden einmal über
\texttt{condor\_master} gestartet und blieben in der verwendeten
Einrichtung dauerhaft aktiv. Vor jedem Aufruf waren zusätzlich zwei
Umgebungsvariablen zu setzen, die nur für die aktuelle
Terminal-Sitzung gelten: \texttt{CONDOR\_CONFIG} für die
HTCondor-Konfiguration und \texttt{PYTHONPATH} für die
Pegasus-Python-Bibliothek. Ohne die zweite Variable brach das
Erzeugungsskript mit einem \texttt{ModuleNotFoundError} ab.

\paragraph{AiiDA:} AiiDA wurde in Version~2.8.0 installiert. Für den in der Dokumentation empfohlenen Betrieb setzt AiiDA zwei Dienste voraus, PostgreSQL als Datenbank-Backend für die Ablage aller Daten- und Prozessknoten und RabbitMQ für die Kommunikation mit dem Daemon~\cite{aiidadocs-install}. Beide Dienste wurden lokal installiert. Über die Installation hinaus waren eine Datenbank und ein Profil anzulegen, vor jedem Lauf muss zusätzlich der Daemon aktiv sein, der die Workflows im Hintergrund ausführt.

Tabelle~\ref{tab:lokal:dienste} fasst die für die lokale
Bereitstellung nötigen Komponenten und Schritte zusammen. Die
Gegenüberstellung bleibt beschreibend, eine Aufwandsrangfolge wird
nicht gebildet, weil sich die Schritte in Art und Umfang
unterscheiden. Die Einordnung aus Einsteigerperspektive folgt in den
nachfolgenden Abschnitten anhand der Vergleichsaspekte aus
Abschnitt~\ref{sec:methodik:qualitativ}.

\begin{table}[ht]
  \caption{Lokale Bereitstellung je System: zusätzliche Komponenten, einmalige Konfiguration und vor der Ausführung erforderliche Schritte}
  \label{tab:lokal:dienste}
  \centering
  \footnotesize
  \renewcommand{\arraystretch}{1.3}
  \begin{tabular}{l l p{2.6cm} p{3.0cm} p{3.4cm}}
    \hline
    \textbf{System} & \textbf{Version} & \textbf{Zusätzliche Komponenten} & \textbf{Einmalige Konfiguration} & \textbf{Vor der Ausführung} \\
    \hline
    Nextflow   & 25.10.4  & keine & keine & keine \\
    StreamFlow & 0.2.0rc2 & keine & keine & keine \\
    Merlin     & 1.13.0   & Redis & \texttt{setuptools}-Downgrade, \texttt{app.yaml} & Redis-Container und Worker starten \\
    Pegasus    & 5.1.2    & HTCondor 23.1.0 & keine & HTCondor-Prozesse aktiv, \texttt{CONDOR\_CONFIG} und \texttt{PYTHONPATH} gesetzt (je Sitzung) \\
    AiiDA      & 2.8.0    & PostgreSQL, RabbitMQ & Datenbank, Profil & Daemon starten \\
    \hline
  \end{tabular}
\end{table}

\section{Pipeline-Pattern}
\label{sec:lokal:pipeline}

Das in Abschnitt~\ref{sec:methodik:versuchsplan} beschriebene
Pipeline-Pattern wurde in allen fünf Systemen mit denselben
Benchmark-Skripten umgesetzt. Im Folgenden wird gegenübergestellt,
wie die fünf Systeme dieses Pattern modellieren und ausführen. Die
vollständigen Workflow-Beschreibungen sind im Repositorium zu
dieser Arbeit abgelegt~\cite{barepo}.

\subsection{Nextflow}
\label{sec:lokal:pipeline:nextflow}

Nextflow beschreibt den Workflow in einer einzigen Datei
\texttt{benchmark.nf} in der systemeigenen Sprache DSL2.

\emph{Ausdruck des Patterns:} Eine separate Abhängigkeitsangabe gibt
es nicht. Die Reihenfolge entsteht dadurch, dass die Ausgabe eines
Prozesses als Eingabe des nachfolgenden Prozesses dient, entsprechend
dem Datenfluss-Modell, auf dem Nextflow
beruht~\cite[S.~317]{ditommaso2017nextflow}.

\begin{lstlisting}[caption={workflow-Block aus benchmark.nf (Nextflow, Pipeline)},label={lst:pipe:nf}]
workflow {
    raw_ch      = GENERATE_INPUT()
    prepared_ch = PREPROCESS(raw_ch)
    result_ch   = COMPUTE(prepared_ch)
    summary_ch  = POSTPROCESS(result_ch)
}
\end{lstlisting}

\emph{Modellierungs- und Konfigurationsaufwand:} Die
Workflow-Beschreibung liegt vollständig in einer Datei, in der
Prozessdefinitionen und Komposition zusammenstehen. Eine Änderung der
Rechenlast betrifft allein den Aufruf des Berechnungsskripts, die
Fassungen für die drei Rechenlasten unterscheiden sich in genau einer
Zeile.

\emph{Sichtbarkeit von Datenfluss und Abhängigkeiten:} Nicht nur die
Reihenfolge, auch der Weg der Dateien ist in der Beschreibung
ablesbar. Die beteiligten Dateien erscheinen in den
\texttt{path}-Deklarationen der Ein- und Ausgaben und eine so
deklarierte Ausgabe wird als Verweis im Arbeitsverzeichnis des
nachfolgenden Prozesses bereitgestellt, sodass die Skripte sie unter
dem angegebenen Namen öffnen. Die Abhängigkeiten entstehen durch die
Übergabe dieser Ausgaben über die Channels an die nachfolgenden
Prozesse.

\emph{Ausführung und Steuerung:} Die Ausführung steuert die
Nextflow-Laufzeit. Ein Prozess wird gestartet, sobald sein
Eingabe-Channel Daten enthält, der Startzeitpunkt jedes Prozesses ergibt sich damit erst zur Laufzeit aus dem Eintreffen der Daten.

\emph{Trennung von Workflow-Logik und Ausführungsumgebung:}
Workflow-Beschreibung und Ausführungsparameter liegen getrennt vor.
Angaben zur Ausführungsumgebung, die Nextflow über sogenannte
Executor konfiguriert, sowie zu den benötigten Ressourcen stehen in
einer eigenen Konfigurationsdatei, sodass ein Wechsel der
Ausführungsumgebung die Workflow-Datei unberührt lässt.


\emph{Nachvollziehbarkeit der Ausführung:} Während des Laufs zeigt die
Ausgabe im Terminal für jeden Prozess den Kennzeichner des
Arbeitsverzeichnisses sowie die Zahl der abgeschlossenen Aufgaben an.
Je Prozess bleibt ein Arbeitsverzeichnis mit den erzeugten Dateien
zurück. Zusätzlich lassen sich eine Trace-Datei, ein Zeitstrahl und
der Workflow-Graph erzeugen~\cite{nextflowdocs-reports}. Der Workflow-Graph ist in
Abbildung~\ref{fig:nf:dag} dargestellt, Zeitstrahlen von
Pipeline-Läufen in den Abbildungen~\ref{fig:nf:timeline:short}
und~\ref{fig:nf:timeline}.

\emph{Verständlichkeit aus Einsteigerperspektive:} Die Abfolge der
vier Schritte ist im \texttt{workflow}-Block kompakt sichtbar. Für das
Verständnis der Beschreibung ist darüber hinaus insbesondere das
Channel-Konzept relevant, über das die Ausgaben eines Prozesses an
den nachfolgenden Prozess weitergegeben werden.

\subsection{StreamFlow}
\label{sec:lokal:pipeline:streamflow}

StreamFlow verteilt die Umsetzung auf drei Dateien:
\texttt{workflow.cwl} enthält die Workflow-Beschreibung,
\texttt{streamflow.yml} die Zuordnung zur Ausführungsumgebung und
\texttt{config.yml} die Eingabewerte und Pfade der Skripte.

\emph{Ausdruck des Patterns:} Eine separate Abhängigkeitsangabe gibt
es nicht. Jeder Schritt benennt unter \texttt{in} seine Eingaben und
eine dieser Eingaben verweist auf die Ausgabe des vorherigen
Schritts. Die Reihenfolge entsteht damit aus den Verweisen zwischen
den Schritten.

\begin{lstlisting}[caption={Verweis auf die Ausgabe des Vorgängers in workflow.cwl (StreamFlow, Pipeline)},label={lst:pipe:sf}]
  preprocess:
    in:
      script:        preprocess_script
      timing_helper: timing_helper
      raw_file:      generate_input/raw_file
    out: [prepared_file, timing_file]
\end{lstlisting}

\emph{Modellierungs- und Konfigurationsaufwand:} Für die
Workflow-Beschreibung und -Konfiguration werden drei Dateien
benötigt: \texttt{workflow.cwl}, \texttt{streamflow.yml} und
\texttt{config.yml}. Eine Änderung der Rechenlast betrifft allein
den Pfad des Berechnungsskripts in der fünfzeiligen
\texttt{config.yml}, deren Fassungen sich in genau einer Zeile
unterscheiden.

\emph{Sichtbarkeit von Datenfluss und Abhängigkeiten:} Die
weitergegebenen Dateien und die daraus entstehenden Abhängigkeiten
sind in den Eingaben der Schritte sichtbar. Jede Eingabe verweist
auf die Workflow-Eingabe oder die Ausgabe eines vorherigen Schritts,
von der sie stammt. So erhält \texttt{preprocess} die Datei
\texttt{raw\_file} über den Verweis
\texttt{generate\_input/raw\_file}. An diesen Verweisen lässt sich
damit sowohl ablesen, welche Dateien zwischen den Schritten
weitergegeben werden, als auch, welche Schritte voneinander
abhängen.

\emph{Ausführung und Steuerung:} Die Ausführung übernimmt die
StreamFlow-Laufzeit. Ein Schritt wird ausgeführt, sobald die
Eingaben vorliegen, auf die er unter \texttt{in} verweist. Für die
lokale Untersuchung wurden alle Schritte einer lokalen
Ausführungsumgebung zugewiesen.

\emph{Trennung von Workflow-Logik und Ausführungsumgebung:} Die
Zuordnung zur Ausführungsumgebung steht vollständig in
\texttt{streamflow.yml} und ist von der fachlichen Beschreibung
getrennt. Diese Aufteilung entspricht dem Entwurfsziel von
StreamFlow, die Workflow-Beschreibung von der Ausführungsumgebung
zu lösen~\cite[S.~1723]{colonnelli2021streamflow}. Ein Wechsel der
Ausführungsumgebung erfordert deshalb keine Änderung an
\texttt{workflow.cwl}, sondern nur an der Angabe des Deployments.


\emph{Nachvollziehbarkeit der Ausführung:} Der Lauf lässt sich
unmittelbar an der Ausgabe im Terminal verfolgen. Sie nennt je Schritt
das Arbeitsverzeichnis und den vollständigen Aufruf mit allen
Argumentpfaden sowie jede einzelne Dateiübertragung mit Quelle und
Ziel. Am Ende gibt StreamFlow die als Workflow-Ausgaben deklarierten
Dateien mit Pfad, Größe und Prüfsumme aus. Aus den gespeicherten
Provenance-Daten lassen sich zudem ein Zeitstrahl des Laufs und ein
Archiv im RO-Crate-Format erzeugen~\cite{streamflowdocs-inspect}. Der
Zeitstrahl eines Pipeline-Laufs ist in
Abbildung~\ref{fig:sf:timeline} dargestellt.

\emph{Verständlichkeit aus Einsteigerperspektive:} Für das
Verständnis der Beschreibung sind die Konzepte der Common Workflow
Language relevant, insbesondere die Typangaben der Ein- und
Ausgaben und die Werkzeugbeschreibung, die je Schritt den Aufruf
des Programms festlegt.

\subsection{Merlin}
\label{sec:lokal:pipeline:merlin}

Merlin beschreibt den Workflow in einer YAML-Spezifikation, deren
Abschnitt \texttt{study} die vier Schritte enthält. Eine zweite Datei
enthält die Verbindungsdaten des Nachrichtenvermittlers.

\emph{Ausdruck des Patterns:} Die Reihenfolge wird ausdrücklich
benannt. Jeder Schritt führt unter \texttt{depends} seine Vorgänger
auf.

\begin{lstlisting}[caption={Abhängigkeit und Datenweitergabe eines Schritts (Merlin, Pipeline)},label={lst:pipe:me}]
  - name: compute_1
    run:
      depends: [preprocess]
      cmd: python3 $(SPECROOT)/scripts/compute_medium.py $(preprocess.workspace)/prepared_input.txt result.txt
\end{lstlisting}

\emph{Modellierungs- und Konfigurationsaufwand:} Für die
Workflow-Beschreibung und -Konfiguration werden zwei Dateien
benötigt: die YAML-Spezifikation mit der Beschreibung des Workflows
und die Datei \texttt{app.yaml} mit den Verbindungsdaten des
Nachrichtenvermittlers. Eine Änderung der Rechenlast betrifft allein
den Aufruf des Berechnungsskripts in der Spezifikation.

\emph{Sichtbarkeit von Datenfluss und Abhängigkeiten:} Die
Abhängigkeiten werden über \texttt{depends} ausdrücklich angegeben.
Welche Datei zwischen zwei Schritten weitergegeben wird, ist dagegen
im jeweiligen Befehl sichtbar. So verweist \texttt{compute} mit
\texttt{\$(preprocess.workspace)/prepared\_input.txt} auf die von
\texttt{preprocess} erzeugte Datei. Abhängigkeit und Datenweitergabe
stehen damit an zwei Stellen desselben Schritts.

\emph{Ausführung und Steuerung:} Die Schritte werden als Aufgaben in
die Warteschlange des Vermittlers eingestellt, aus der die
Celery-Worker sie entgegennehmen und
abarbeiten~\cite[S.~258]{peterson2022merlin}\cite{merlindocs-celery}.
Ein Schritt wird dabei erst zur Ausführung freigegeben, wenn seine
unter \texttt{depends} angegebenen Vorgänger abgeschlossen sind.

\emph{Trennung von Workflow-Logik und Ausführungsumgebung:} Die
Trennung beruht auf dem Erzeuger-Verbraucher-Modell. Die
Spezifikation beschreibt, welche Schritte auszuführen sind, während
die Worker die Aufgaben aus der Warteschlange abholen, unabhängig
davon, wo sie laufen. Die Workflow-Beschreibung ist damit von der
ausführenden Ressource entkoppelt~\cite[S.~261]{peterson2022merlin}. Die
Spezifikation benennt zwar, welche Schritte ein Worker übernimmt,
nicht aber, auf welcher Ressource dieser Worker läuft.


\emph{Nachvollziehbarkeit der Ausführung:} Jeder Schritt läuft in
einem eigenen Arbeitsverzeichnis, sodass sich Ergebnisse und
Protokolle eindeutig einem Schritt zuordnen lassen. Während der
Ausführung zeigen die Ausgaben der Celery-Worker fortlaufend an,
welche Aufgaben angenommen und abgeschlossen werden. Über die
Kommandozeile lassen sich zudem der Zustand der Warteschlangen und
Worker sowie der Fortschritt eines Laufs abfragen~\cite{merlindocs-commands}.
Eine Funktion zur Ausgabe des Workflow-Graphen wurde in der
verwendeten Version~1.13.0 nicht gefunden.

\emph{Verständlichkeit aus Einsteigerperspektive:} Die Vorgänger
jedes Schritts sind über \texttt{depends} ausdrücklich bezeichnet.
Für das Verständnis der Beschreibung sind darüber hinaus die
Workspace-Verweise und das Zusammenspiel von Nachrichtenvermittler
und Workern relevant.

\subsection{Pegasus}
\label{sec:lokal:pipeline:pegasus}

Der Workflow wird in \texttt{build\_workflow.py} über die
Python-Schnittstelle von Pegasus aufgebaut.

\emph{Ausdruck des Patterns:} Die Reihenfolge ergibt sich aus den
Dateien. Jeder Job benennt über \texttt{add\_inputs} und
\texttt{add\_outputs}, welche Dateien er liest und schreibt und
Pegasus leitet daraus die Kanten des Graphen ab.

\begin{lstlisting}[language=Python,caption={Ein- und Ausgaben eines Jobs (Pegasus, Pipeline)},label={lst:pipe:pg}]
compute_job = Job("compute")
compute_job.add_args("prepared_input.txt", "result.txt")
compute_job.add_inputs(prepared_file)
compute_job.add_outputs(result_file, stage_out=False, register_replica=False)
\end{lstlisting}

\emph{Modellierungs- und Konfigurationsaufwand:} Für die
Workflow-Beschreibung und -Konfiguration wurden vier Dateien
manuell erstellt: \texttt{build\_workflow.py}, \texttt{sites.yml},
\texttt{pegasus.properties} und \texttt{plan.sh}. Das
Erzeugungsskript erzeugt die Workflow-Datei und den
Transformationskatalog. Eine Änderung der Rechenlast betrifft allein
den Pfad des Berechnungsskripts im Erzeugungsskript. Die erzeugten
Workflow-Dateien der drei Rechenlasten sind identisch.

\emph{Sichtbarkeit von Datenfluss und Abhängigkeiten:} Die
weitergegebenen Dateien werden als eigene \texttt{File}-Objekte
angelegt und bei den Jobs als Ein- und Ausgaben angegeben. So
erzeugt \texttt{preprocess} die Datei
\texttt{prepared\_input.txt}, die \texttt{compute} anschließend als
Eingabe verwendet. Die Abhängigkeiten zwischen den Jobs entstehen
damit aus diesen Ein- und Ausgaben.

\emph{Ausführung und Steuerung:} Pegasus trennt Planung und
Ausführung. Der Planer übersetzt die abstrakte Beschreibung in einen
konkreten Graphen, der anschließend an HTCondor zur Ausführung
übergeben wird~\cite[S.~20]{deelman2015pegasus}. Welche Jobs mit welchen
Abhängigkeiten auszuführen sind, steht damit bereits vor Beginn der
Ausführung fest.

\emph{Trennung von Workflow-Logik und Ausführungsumgebung:} Die
Trennung beruht auf der abstrakten Workflow-Beschreibung. Der
Workflow wird ressourcenunabhängig formuliert, während der
Transformations- und der Site-Katalog die Verbindung zu den
Programmen und den verfügbaren Rechenumgebungen herstellen und der
Planer daraus den ausführbaren Workflow
erzeugt~\cite[S.~18--20]{deelman2015pegasus}.


\emph{Nachvollziehbarkeit der Ausführung:} Die Jobs eines Laufs
werden in einem gemeinsamen Arbeitsverzeichnis ausgeführt, ein
eigenes Verzeichnis je Job entsteht nicht. Während des Laufs lässt
sich der Fortschritt über die Statusabfrage verfolgen, die die
laufenden Jobs einzeln benennt und den Anteil der abgeschlossenen
Jobs ausweist. Bei jedem Lauf entstanden
im Submit-Verzeichnis neben den Protokollen automatisch drei
Abbildungen des Workflow-Graphen, eine abstrakte Darstellung, eine
Darstellung mit den beteiligten Dateien und eine des konkreten,
geplanten Graphen. Die drei Darstellungen des
Pipeline-Workflows sind in den Abbildungen~\ref{fig:pg:dag},
\ref{fig:pg:dag:files} und~\ref{fig:pg:dag:concrete} dargestellt.

\emph{Verständlichkeit aus Einsteigerperspektive:} Die Beschreibung
erfolgt in Python. Für das Verständnis der Umsetzung sind darüber
hinaus die Kataloge, der vorgelagerte
Planungsschritt und die Rolle von HTCondor relevant.

\subsection{AiiDA}
\label{sec:lokal:pipeline:aiida}

In AiiDA wird der Workflow als Python-Klasse beschrieben. Eine
WorkChain legt den Ablauf fest, ein allgemein gehaltener CalcJob
führt je Schritt eines der Benchmark-Skripte aus, ein Submit-Skript
übergibt den Lauf an den Daemon.

\emph{Ausdruck des Patterns:} Die Reihenfolge wird ausdrücklich als
Liste von Methoden in \texttt{spec.outline()} festgelegt. Die
abschließende Methode \texttt{results} reicht keinen weiteren
CalcJob ein, sie veröffentlicht die Ausgaben der vier fachlichen
Schritte als Ausgaben der WorkChain.

\begin{lstlisting}[language=Python,caption={Reihenfolge in spec.outline() (AiiDA, Pipeline)},label={lst:pipe:ai}]
spec.outline(
    cls.submit_generate,
    cls.submit_preprocess,
    cls.submit_compute,
    cls.submit_postprocess,
    cls.results,
)
\end{lstlisting}

\emph{Modellierungs- und Konfigurationsaufwand:} Die Umsetzung
umfasst WorkChain, CalcJob und Submit-Skript, wobei der CalcJob
allgemein formuliert ist und in beiden Mustern unverändert verwendet
wird. Eine eigene Fassung je Rechenlast entfällt, die Rechenlast wird
dem Submit-Skript als Argument übergeben und die WorkChain wählt
daraus das Berechnungsskript aus.

\emph{Sichtbarkeit von Datenfluss und Abhängigkeiten:} Die
weitergegebenen Dateien werden in der WorkChain ausdrücklich den
aufeinanderfolgenden Schritten zugeordnet. So wird die Ausgabe eines
CalcJobs als Eingabe des nachfolgenden CalcJobs verwendet. Welche
Dateien zwischen den Schritten weitergegeben werden, ist damit in der
WorkChain sichtbar. Die Abhängigkeiten ergeben sich aus der
in \texttt{spec.outline()} festgelegten Reihenfolge der Schritte.

\emph{Ausführung und Steuerung:} Das Submit-Skript übergibt die
WorkChain an den Daemon, der die Schritte der Outline abarbeitet und
für jeden Schritt einen CalcJob einreicht. Der jeweils nachfolgende
Schritt wird erst ausgeführt, wenn der vorherige CalcJob erfolgreich
abgeschlossen wurde.

\emph{Trennung von Workflow-Logik und Ausführungsumgebung:} Die
Ausführungsumgebung wird als Eingabe behandelt. Rechner und Code sind
in der Datenbank als eigene Objekte abgelegt und werden vom
Submit-Skript geladen. Ein Wechsel der Ausführungsumgebung erfordert
damit keine Änderung der WorkChain, sondern die Verwendung
entsprechend eingerichteter Rechner- und Code-Objekte.


\emph{Nachvollziehbarkeit der Ausführung:} Während des Laufs lässt
sich der Fortschritt über die Statusabfrage verfolgen, die die
laufenden Prozesse einzeln benennt und ihren Zustand ausweist. Unabhängig davon hält AiiDA jeden Schritt bereits während des Laufs in einem dauerhaften
Provenance-Graphen fest. Darin sind sowohl die Prozesse als auch die
Daten eigene Knoten und die Kanten geben an, welche Eingaben ein
Prozess verwendet hat und welche Ausgaben aus ihm hervorgegangen
sind~\cite[S.~7\,f.]{huber2020aiida}. Dabei werden zwei Zusammenhänge unterschieden. Zum einen wird festgehalten, wie Daten aus vorherigen Daten und Berechnungsschritten entstanden sind. Zum anderen wird dokumentiert, welcher Workflow die einzelnen Berechnungen ausgelöst hat. In der vorliegenden Umsetzung erscheint deshalb die WorkChain als
eigener Knoten, der mit den vier von ihr eingereichten CalcJobs
verbunden ist. Der gespeicherte Ablauf lässt sich im
Nachhinein abfragen und als Graph ausgeben~\cite{aiidadocs-graph}. Der
Provenance-Graph des Pipeline-Laufs ist in
Abbildung~\ref{fig:ai:prov} dargestellt.

\emph{Verständlichkeit aus Einsteigerperspektive:} Die Beschreibung
erfolgt in Python. Für das Verständnis der Umsetzung sind darüber hinaus die
Datenobjekte von AiiDA, das Zusammenspiel von WorkChain und CalcJob,
der Kontext für die Weitergabe von Ergebnissen sowie die Einrichtung
von Rechner und Code relevant.

\subsection{Vergleichende Beobachtungen zum Pipeline-Pattern}
\label{sec:lokal:pipeline:zusammenfassung}

Tabelle~\ref{tab:pipe:vergleich} stellt die Beobachtungen aus den
fünf Systemabschnitten entlang der Vergleichsaspekte gegenüber. Als
Artefakte gezählt werden dabei die manuell erstellten Dateien, die
die Workflow-Beschreibung und ihre Konfiguration ausmachen,
automatisch erzeugte Dateien und die unveränderten Benchmark-Skripte
bleiben unberücksichtigt. Bei AiiDA treten zu den drei Dateien
Rechner und Code hinzu, die einmalig als Objekte in der Datenbank
eingerichtet werden.

\begin{table}[H]
  \caption{Umsetzung des Pipeline-Patterns im Überblick}
  \label{tab:pipe:vergleich}
  \centering
  \scriptsize
  \setlength{\tabcolsep}{3pt}
  \renewcommand{\arraystretch}{1.4}
  \begin{tabular}{>{\raggedright\arraybackslash}p{2.7cm} >{\raggedright\arraybackslash}p{1.9cm} >{\raggedright\arraybackslash}p{1.9cm} >{\raggedright\arraybackslash}p{1.9cm} >{\raggedright\arraybackslash}p{1.9cm} >{\raggedright\arraybackslash}p{1.9cm}}
    \hline
     & \textbf{Nextflow} & \textbf{Stream\-Flow} & \textbf{Merlin} & \textbf{Pegasus} & \textbf{AiiDA} \\
    \hline
    Ausdruck des Patterns & Übergabe der Prozess\-ausgaben & Verweis auf Vorgänger\-ausgabe & Angabe unter \texttt{depends} & deklarierte Ein- und Ausgaben & Liste in \texttt{spec.outline()} \\
    Modellierung und Konfiguration & eine Datei & drei Dateien & zwei Dateien & vier Dateien & drei Dateien, dazu Rechner und Code \\
    Datenfluss und Abhängigkeiten & \texttt{path}-Angaben und Channels & Verweise auf Schritt\-ausgaben & Pfad im Befehl & \texttt{File}-Objekte der Jobs & Zuordnung in der WorkChain \\
    Ausführung und Steuerung & eigene Laufzeit & eigene Laufzeit & Celery-Worker & HTCondor nach Planung & Daemon \\
    Trennung Logik und Umgebung & eigene Konfigu\-rationsdatei & eigene Deployment-Datei & Warteschlange trennt Schritt und Ressource & abstrakte Beschreibung, Kataloge, Planer & Rechner und Code in der Datenbank \\
    Nachvollziehbarkeit & Fortschritts\-anzeige, Trace-Datei, Zeitstrahl, Graph & Terminal\-ausgabe, Zeitstrahl, RO-Crate & Worker-Ausgaben, Status\-abfragen & Status\-abfragen, drei Graph\-abbildungen & Status\-abfragen, Provenance-Graph \\
    \hline
  \end{tabular}
\end{table}

Beim Ausdruck des Patterns zeigen sich zwei Ansätze. Nextflow,
StreamFlow und Pegasus leiten die Reihenfolge der Schritte aus den
Datenbeziehungen ab, während Merlin und AiiDA die Abfolge
ausdrücklich beschreiben.

Beim Modellierungs- und Konfigurationsaufwand reicht die Umsetzung
von einer Datei bei Nextflow über mehrere Beschreibungs- und
Konfigurationsdateien bei StreamFlow, Merlin und Pegasus bis zu drei
Python-Bausteinen bei AiiDA. Eine Änderung der Rechenlast betrifft
bei Nextflow, StreamFlow, Merlin und Pegasus jeweils nur eine Stelle,
jedoch in unterschiedlichen Dateien. AiiDA behandelt die Rechenlast
stattdessen als Aufrufargument.

Auch Datenfluss und Abhängigkeiten werden unterschiedlich sichtbar.
Nextflow und StreamFlow deklarieren die weitergegebenen Dateien an
den Schritten, Pegasus führt sie als eigene Objekte, bei Merlin
stehen Abhängigkeit und Dateipfad getrennt und AiiDA ordnet die
Ausgaben eines CalcJobs dem nachfolgenden CalcJob in der WorkChain
zu.

Bei der Ausführung steuern Nextflow und StreamFlow die Schritte über
ihre eigene Laufzeit. Merlin übergibt Aufgaben an Celery-Worker,
Pegasus den zuvor geplanten Workflow an HTCondor und AiiDA die
WorkChain an den Daemon. Bei Pegasus steht der auszuführende Graph
damit bereits vor der Ausführung fest, während bei Nextflow der Start
eines Prozesses durch eintreffende Daten ausgelöst wird.

Bei der Trennung von Logik und Umgebung greifen unterschiedliche
Mechanismen. Nextflow und StreamFlow lagern die Umgebungsangaben in
eigene Dateien aus, Merlin entkoppelt Beschreibung und Ressource über
die Aufgabenwarteschlange, Pegasus formuliert den Workflow abstrakt
und stellt die Verbindung zur Umgebung über Kataloge und den Planer
her und AiiDA hält Rechner und Code als Objekte in der Datenbank.

Bei der Nachvollziehbarkeit bieten alle fünf Systeme Möglichkeiten,
den Fortschritt während eines Laufs zu verfolgen, unterscheiden sich
aber in den verbleibenden Artefakten. Nextflow stellt unter anderem
Trace-Datei, Zeitstrahl und Workflow-Graph bereit, StreamFlow einen
Zeitstrahl und Provenance-Daten, Merlin vor allem Arbeitsverzeichnisse
und Worker-Protokolle und Pegasus Protokolle sowie Darstellungen des
abstrakten und geplanten Graphen. AiiDA unterscheidet sich durch den
dauerhaft gespeicherten Provenance-Graphen, der Prozesse, Daten und
deren Entstehungszusammenhänge festhält.

Aus Einsteigerperspektive unterscheiden sich die Umsetzungen vor
allem darin, wie viele systemeigene Konzepte vor dem ersten Lauf
verstanden sein müssen. Bei Nextflow ist insbesondere das Channel-Modell relevant, bei StreamFlow die Common Workflow Language, bei Merlin das
Zusammenspiel von Nachrichtenvermittler und Workern und bei Pegasus
Kataloge, Planung und HTCondor. AiiDA erfordert das Verständnis von
Datenobjekten, WorkChain und CalcJob sowie der Einrichtung von
Rechner und Code. Diese Einordnung beruht auf der erstmaligen
Umsetzung und ersetzt keine Nutzerstudie. Wie sich die beschriebenen Unterschiede im
Laufzeitverhalten niederschlagen, untersucht
Abschnitt~\ref{sec:lokal:messungen}.

\section{Scatter-Gather-Pattern}
\label{sec:lokal:scattergather}

Das in Abschnitt~\ref{sec:methodik:versuchsplan} beschriebene
Scatter-Gather-Pattern wurde ebenfalls in allen fünf Systemen mit
denselben Benchmark-Skripten umgesetzt. Untersucht wurden Läufe mit
einem, zwei und vier Teilstücken. Im Folgenden wird
gegenübergestellt, wie die fünf Systeme dieses Pattern modellieren
und ausführen. Die vollständigen Workflow-Beschreibungen sind im
Repositorium zu dieser Arbeit abgelegt~\cite{barepo}.

\subsection{Nextflow}
\label{sec:lokal:sg:nextflow}

Die Beschreibung folgt dem Aufbau des Pipeline-Workflows und liegt
wieder vollständig in einer Datei, die Zahl der Teilstücke steht als
Parameter am Dateianfang.

\emph{Ausdruck des Patterns:} Die parallele Phase wird über die
Operatoren \texttt{flatten} und \texttt{collect} auf den Channels
ausgedrückt. \texttt{SPLIT} erzeugt die Teilstücke als gemeinsame
Ausgabe, \texttt{flatten} gibt sie als einzelne Elemente weiter,
sodass \texttt{COMPUTE} für jedes Element gestartet wird.
\texttt{collect} sammelt anschließend die Teilergebnisse zu einer
gemeinsamen Eingabe für \texttt{AGGREGATE}.

\begin{lstlisting}[caption={Fan-out und Fan-in im workflow-Block (Nextflow, Scatter-Gather)},label={lst:sg:nf}]
    chunks_ch = SPLIT(prepared_ch)
        .flatten()
    results_ch = COMPUTE(chunks_ch)
        .collect()
    aggregated_ch = AGGREGATE(results_ch)
\end{lstlisting}

\emph{Modellierungs- und Konfigurationsaufwand:} Die Zahl der
Teilstücke ist ein Parameter, die Fassungen für einen, zwei und vier
Teilstücke unterscheiden sich in genau einer Zeile. Der Wechsel der
Rechenlast betrifft wie beim Pipeline-Pattern allein den Aufruf des
Berechnungsskripts, ebenfalls eine Zeile.

\emph{Sichtbarkeit von Datenfluss und Abhängigkeiten:} Der Weg der
Dateien durch die parallele Phase ist in der Beschreibung ablesbar.
\texttt{SPLIT} deklariert seine Ausgabe als Dateimuster, an dem
sichtbar ist, dass mehrere Dateien entstehen. Nach
\texttt{flatten} erhält jede Instanz von \texttt{COMPUTE} genau ein
Teilstück als Eingabe, während \texttt{AGGREGATE} nach
\texttt{collect} die Ergebnisdateien aller Instanzen erhält. Die
Abhängigkeiten entstehen wie beim Pipeline-Pattern durch die
Übergabe dieser Ausgaben über die Channels an die nachfolgenden
Prozesse.

\emph{Ausführung und Steuerung:} Für jedes Element des
Eingabe-Channels startet die Nextflow-Laufzeit eine eigene Instanz
von \texttt{COMPUTE}, die Instanzen laufen nebeneinander.
\texttt{AGGREGATE} startet erst, wenn \texttt{collect} die
Ergebnisse aller Instanzen zusammengeführt hat.

\emph{Trennung von Workflow-Logik und Ausführungsumgebung:} Wie beim
Pipeline-Pattern steht die Ausführungsumgebung in der eigenen
Konfigurationsdatei, die Workflow-Datei bleibt bei einem Wechsel
unverändert.


\emph{Nachvollziehbarkeit der Ausführung:} Die Fortschrittsanzeige
zählt die parallelen Instanzen von \texttt{COMPUTE} als Aufgaben
desselben Prozesses. Workflow-Graph und Zeitstrahl stehen wie beim
Pipeline-Pattern zur Verfügung, der Graph des
Scatter-Gather-Workflows und der Zeitstrahl eines Laufs mit vier
Teilstücken sind in den Abbildungen~\ref{fig:nf:dag:sg}
und~\ref{fig:nf:timeline:sg} dargestellt.

\emph{Verständlichkeit aus Einsteigerperspektive:} Zusätzlich zum
Channel-Konzept sind die beiden Operatoren \texttt{flatten} und
\texttt{collect} zu verstehen, die mehrere ausgegebene Elemente
einzeln weitergeben beziehungsweise wieder zu einer gemeinsamen
Eingabe sammeln.

\subsection{StreamFlow}
\label{sec:lokal:sg:streamflow}

Die Umsetzung des Scatter-Gather-Patterns verwendet wie beim
Pipeline-Pattern \texttt{workflow.cwl}, \texttt{streamflow.yml}
und \texttt{config.yml}.

\emph{Ausdruck des Patterns:} Die Aufteilung wird über das
\texttt{scatter}-Attribut des \texttt{compute}-Schritts ausgedrückt.
Der Schritt erhält das von \texttt{split} erzeugte Feld der
Teilstücke und wird für jedes Element einzeln ausgeführt.
Für die Zusammenführung verweist \texttt{aggregate} mit
\texttt{result\_files: compute/result\_file} auf die Ausgaben aller
Instanzen und erhält sie als eine Liste.

\begin{lstlisting}[caption={Fan-out über scatter am compute-Schritt (StreamFlow, Scatter-Gather)},label={lst:sg:sf}]
  compute:
    in:
      chunk_file:  split/chunk_files
      output_name: output_names
    scatter: [chunk_file, output_name]
    scatterMethod: dotproduct
\end{lstlisting}

\emph{Modellierungs- und Konfigurationsaufwand:} Wie beim
Pipeline-Pattern werden drei Dateien für Beschreibung und
Konfiguration verwendet. Die \texttt{workflow.cwl} bleibt über alle
Varianten unverändert. Eine Änderung der Teilstückzahl betrifft allein
die \texttt{config.yml}. Dort ändern sich der Eingabewert und die
Liste der Ergebnisnamen, die einen Eintrag je Teilstück enthält. Der
Wechsel der Rechenlast betrifft ebenfalls nur die
\texttt{config.yml}, dort ändert sich der Pfad des
Berechnungsskripts.

\emph{Sichtbarkeit von Datenfluss und Abhängigkeiten:} Wie beim
Pipeline-Pattern sind die weitergegebenen Dateien und ihre Herkunft
in den Eingaben der Schritte sichtbar. Zusätzlich zeigen die
Typangaben, ob eine einzelne Datei oder ein Feld von Dateien
weitergegeben wird. \texttt{split} erzeugt ein Feld von Teilstücken,
aus dem jede Ausführung von \texttt{compute} ein Teilstück erhält.
\texttt{aggregate} verweist anschließend auf die Ausgaben von
\texttt{compute} und erhält das Feld aller Ergebnisdateien. Die
Abhängigkeiten zwischen den Schritten ergeben sich damit wie beim
Pipeline-Pattern aus diesen Verweisen auf die jeweiligen Ausgaben.

\emph{Ausführung und Steuerung:} Die Laufzeit startet für jedes
Element eine eigene Instanz des \texttt{compute}-Schritts, die
Instanzen laufen parallel in der zugewiesenen Umgebung.
\texttt{aggregate} startet erst, wenn die Ergebnisse aller Instanzen
vorliegen.

\emph{Trennung von Workflow-Logik und Ausführungsumgebung:} Wie beim
Pipeline-Pattern steht die Zuordnung zur Ausführungsumgebung in
\texttt{streamflow.yml}.


\emph{Nachvollziehbarkeit der Ausführung:} Der Lauf lässt sich wie
beim Pipeline-Pattern an der Terminalausgabe verfolgen, Zeitstrahl
und RO-Crate-Archiv stehen ebenfalls zur Verfügung. Der Zeitstrahl
eines Scatter-Gather-Laufs ist in
Abbildung~\ref{fig:sf:timeline:sg} dargestellt.

\emph{Verständlichkeit aus Einsteigerperspektive:} Zusätzlich zu den
Konzepten des Pipeline-Patterns ist das \texttt{scatter}-Attribut zu
verstehen sowie die Unterscheidung, ob eine Ein- oder Ausgabe eine
einzelne Datei oder eine Sammlung von Dateien bezeichnet.

\subsection{Merlin}
\label{sec:lokal:sg:merlin}

Merlin beschreibt auch das Scatter-Gather-Pattern in einer
YAML-Spezifikation.

\emph{Ausdruck des Patterns:} Die parallele Verarbeitung wird über
den Parameter \texttt{CHUNK} ausgedrückt, dessen Werteliste einen
Eintrag je Teilstück enthält. Der einmal beschriebene
\texttt{compute}-Schritt wird dadurch für jeden Parameterwert
instanziiert. Jede Instanz hängt von \texttt{split} ab. Für den
anschließenden Fan-in verweist \texttt{aggregate} unter
\texttt{depends} mit \texttt{compute\_*} auf alle hervorgegangenen Compute-Instanzen.

\begin{lstlisting}[caption={Parametrisierter Verarbeitungsschritt und Fan-in (Merlin, Scatter-Gather)},label={lst:sg:me}]
  - name: compute
    run:
      depends: [split]
      cmd: python3 $(SPECROOT)/scripts/compute_medium.py $(split.workspace)/chunk_$(CHUNK).txt result_$(CHUNK).txt

  - name: aggregate
    run:
      depends: [compute_*]
\end{lstlisting}

\emph{Modellierungs- und Konfigurationsaufwand:} Die Zahl der
Teilstücke steht an zwei Stellen der Spezifikation, in der
Werteliste des Parameters und im Argument des Zerlegeschritts. Der
Wechsel der Rechenlast betrifft wie beim Pipeline-Pattern allein den
Aufruf des Berechnungsskripts.

\emph{Sichtbarkeit von Datenfluss und Abhängigkeiten:} Die
Abhängigkeiten werden wie beim Pipeline-Pattern über
\texttt{depends} ausdrücklich angegeben. Welche Dateien
weitergegeben werden, ist in den Befehlen sichtbar. Jede
\texttt{compute}-Instanz liest über den jeweiligen Wert von
\texttt{CHUNK} eine Datei aus dem Arbeitsverzeichnis von
\texttt{split} und erzeugt die entsprechende Ergebnisdatei.
\texttt{aggregate} verweist anschließend auf die Ergebnisdateien
aller Compute-Instanzen. Abhängigkeiten und Datenweitergabe stehen
damit auch im Scatter-Gather-Pattern an getrennten Stellen der
Beschreibung.

\emph{Ausführung und Steuerung:} Die aus der Parametrisierung
hervorgegangenen Compute-Instanzen werden wie beim Pipeline-Pattern
als Aufgaben in die Warteschlange eingestellt und von den
Celery-Workern ausgeführt. Eine Compute-Instanz wird nach Abschluss
von \texttt{split} zur Ausführung freigegeben. \texttt{aggregate}
wird erst freigegeben, wenn alle unter
\texttt{depends: [compute\_*]} erfassten Compute-Instanzen
abgeschlossen sind. In der lokalen Untersuchung konnte der Worker
laut Spezifikation bis zu vier Aufgaben gleichzeitig bearbeiten.

\emph{Trennung von Workflow-Logik und Ausführungsumgebung:} Die
Trennung entspricht dem Pipeline-Pattern. Die Compute-Instanzen
werden als Aufgaben über die Warteschlange an die Worker übergeben,
ohne dass ihre ausführende Ressource in den einzelnen
Workflow-Schritten festgelegt wird. Die Anzahl gleichzeitig
bearbeitbarer Aufgaben wird für den Worker im Abschnitt
\texttt{resources} konfiguriert.

\emph{Nachvollziehbarkeit der Ausführung:} Die einzelnen
Compute-Instanzen erscheinen in den Ausgaben des Workers mit eigenem
Namen, Arbeitsverzeichnis und Abschlussmeldung und hinterlassen
jeweils ein eigenes Arbeitsverzeichnis. Statusabfragen über die
Kommandozeile stehen wie beim Pipeline-Pattern zur Verfügung, eine
Graphausgabe weiterhin nicht.

\emph{Verständlichkeit aus Einsteigerperspektive:} Zusätzlich zu den
Konzepten des Pipeline-Patterns ist das Parameter-Konzept zu
verstehen. Zu beachten ist außerdem, dass die Werteliste und die im Zerlegeschritt angegebene Zahl der Teilstücke zueinander passen müssen.

\subsection{Pegasus}
\label{sec:lokal:sg:pegasus}

Beim Scatter-Gather-Pattern wird der Workflow ebenfalls über das
Erzeugungsskript \texttt{build\_workflow.py} aufgebaut.

\emph{Ausdruck des Patterns:} Die parallele Phase entsteht im
Erzeugungsskript. Eine Schleife legt je Teilstück einen eigenen
\texttt{compute}-Job an, der sein Teilstück als Eingabe und die
zugehörige Ergebnisdatei als Ausgabe deklariert, die Kanten leitet
Pegasus wie beim Pipeline-Pattern aus den Dateien ab. Für die
Zusammenführung deklariert der \texttt{aggregate}-Job die
Ergebnisdateien aller Teilstücke als Eingaben. In der erzeugten
Workflow-Datei stehen die Verarbeitungsjobs einzeln, der Umfang der
parallelen Phase liegt damit bereits in der abstrakten Beschreibung
fest.

\begin{lstlisting}[language=Python,caption={Erzeugung der compute-Jobs je Teilstück (Pegasus, Scatter-Gather)},label={lst:sg:pg}]
for index, (chunk_file, result_file) in enumerate(zip(chunk_files, result_files), start=1):
    compute_job = Job("compute")
    compute_job.add_args(f"chunk_{index}.txt", f"result_{index}.txt")
    compute_job.add_inputs(chunk_file)
    compute_job.add_outputs(result_file, stage_out=False, register_replica=False)
    compute_jobs.append(compute_job)
\end{lstlisting}

\emph{Modellierungs- und Konfigurationsaufwand:} Die vier für das
Pipeline-Pattern benötigten Dateien bleiben bestehen. Das
Erzeugungsskript ist für alle Teilstückzahlen einer Rechenlast
dasselbe, die Teilstückzahl wird ihm als Aufrufargument übergeben.
Die erzeugte Workflow-Datei wächst mit der Teilstückzahl, da je
Teilstück ein Verarbeitungsjob und die zugehörigen Dateiobjekte
hinzukommen. Eine Änderung der Rechenlast betrifft wie beim
Pipeline-Pattern allein den Pfad des Berechnungsskripts im
Erzeugungsskript.

\emph{Sichtbarkeit von Datenfluss und Abhängigkeiten:} Wie beim
Pipeline-Pattern werden die weitergegebenen Dateien als eigene
\texttt{File}-Objekte angelegt und bei den Jobs als Ein- und
Ausgaben angegeben. Die parallele Phase ist daran ablesbar, dass
\texttt{split} alle Teilstücke als Ausgaben deklariert, jeder
\texttt{compute}-Job genau eines dieser Teilstücke als Eingabe
verwendet und \texttt{aggregate} die Ergebnisdateien aller
\texttt{compute}-Jobs als Eingaben erhält. Die Abhängigkeiten
zwischen den Jobs entstehen damit aus diesen Ein- und Ausgaben.

\emph{Ausführung und Steuerung:} Planung und Ausführung bleiben
getrennt. Bei vier Teilstücken enthält die abstrakte Beschreibung
neun fachliche Jobs. Der Planer übersetzt sie in den ausführbaren Graphen, den HTCondor abarbeitet. Dabei ergänzt er Jobs für das Anlegen von Verzeichnissen, das Sichern der Ergebnisse und das Aufräumen.

\emph{Trennung von Workflow-Logik und Ausführungsumgebung:} Die
Trennung über die abstrakte Beschreibung und die Kataloge bleibt wie
beim Pipeline-Pattern bestehen, ebenso die Einschränkung, dass das
Erzeugungsskript die Zielsite benennt.


\emph{Nachvollziehbarkeit der Ausführung:} Die drei Abbildungen des
Workflow-Graphen entstehen wie beim Pipeline-Pattern bei jedem Lauf.
In der abstrakten Darstellung sind die Verarbeitungsjobs einzeln
sichtbar, im geplanten Graphen kommen die vom Planer ergänzten Jobs
hinzu. Die drei Darstellungen des Scatter-Gather-Workflows sind in
den Abbildungen~\ref{fig:pg:dag:sg}, \ref{fig:pg:dag:sg:files}
und~\ref{fig:pg:dag:sg:concrete} dargestellt.

\emph{Verständlichkeit aus Einsteigerperspektive:} Zusätzlich zu den
Konzepten des Pipeline-Patterns ist zu verstehen, dass die parallele
Phase nicht über ein Sprachmittel des Systems ausgedrückt, sondern
im Erzeugungsskript ausgeschrieben wird. In der erzeugten
Beschreibung ist sie nur daran erkennbar, dass die
Verarbeitungsjobs keine Dateien voneinander lesen.

\subsection{AiiDA}
\label{sec:lokal:sg:aiida}

Der Aufbau aus WorkChain, CalcJob und Submit-Skript bleibt bestehen.
Die Outline wächst um das Zerlegen und das Zusammenführen, die
parallele Verarbeitung bleibt darin ein einzelner Eintrag.
Teilstückzahl und Rechenlast sind Eingaben der WorkChain.

\emph{Ausdruck des Patterns:} Die parallele Phase entsteht in einem
einzigen Outline-Schritt. Eine Schleife reicht je Teilstück einen
eigenen CalcJob ein, der sein Teilstück verarbeitet und
\texttt{ToContext} lässt die WorkChain auf den Abschluss aller
eingereichten CalcJobs warten, bevor der nächste Outline-Schritt
beginnt. Dort sammelt \texttt{submit\_aggregate} die Ergebnisordner
aller Instanzen ein und reicht einen einzelnen CalcJob für das
Zusammenführen ein.

\begin{lstlisting}[language=Python,caption={Fan-out in submit\_compute (AiiDA, Scatter-Gather)},label={lst:sg:ai}]
    def submit_compute(self):
        futures = {}
        for index in range(1, self.inputs.num_chunks.value + 1):
            future = self.submit_task(
                label=f"compute_{index}",
                arguments=[f"chunk_{index}.txt", f"result_{index}.txt"],
            )
            futures[f"compute_{index}"] = future
        return ToContext(**futures)
\end{lstlisting}

\emph{Modellierungs- und Konfigurationsaufwand:} Eine eigene Fassung
je Teilstückzahl oder Rechenlast entfällt. Beide werden dem
Submit-Skript als Argumente übergeben und gehen als Eingaben in die
WorkChain ein. Die Rechenlast wählt wie beim Pipeline-Pattern das
Berechnungsskript aus.

\emph{Sichtbarkeit von Datenfluss und Abhängigkeiten:} Die Weitergabe
ist wie beim Pipeline-Pattern an drei Stellen beschrieben, als
Quellordner, als Zuordnung und als Abrufliste, nun jedoch je Instanz.
Jede Verarbeitungsinstanz erhält den Ordner des Zerlegeschritts als
Quelle mit der Zuordnung genau ihres Teilstücks. Beim Zusammenführen
erhält der \texttt{aggregate}-CalcJob die Ordner aller Instanzen als
Quellen mit einer Zuordnung je Ergebnisdatei.

\emph{Ausführung und Steuerung:} Die WorkChain wird weiterhin über
den Daemon ausgeführt. Die Ausführung übernehmen dabei mehrere parallel laufende Runner, die Aufgaben unabhängig voneinander verarbeiten. Da die
Verarbeitungsinstanzen gemeinsam eingereicht werden, können ihre
CalcJobs gleichzeitig verarbeitet werden. Durch \texttt{ToContext}
setzt die WorkChain erst fort, wenn alle eingereichten CalcJobs
beendet sind.


\emph{Trennung von Workflow-Logik und Ausführungsumgebung:} Die
Trennung entspricht dem Pipeline-Pattern: Rechner und Code sind in
der AiiDA-Datenbank abgelegt und werden der WorkChain als Eingaben
übergeben.


\emph{Nachvollziehbarkeit der Ausführung:} Der Provenance-Graph
enthält die Verarbeitungsinstanzen als eigene Prozessknoten, die
sämtlich mit der WorkChain verbunden sind. Die Kanten zeigen, aus
welchem Ordnerknoten jede Instanz ihr Teilstück bezogen hat und
welche Ergebnisse in das Zusammenführen eingegangen sind. Der
Provenance-Graph eines Scatter-Gather-Laufs ist in
Abbildung~\ref{fig:ai:prov:sg} dargestellt.

\emph{Verständlichkeit aus Einsteigerperspektive:} Zusätzlich zu den
Konzepten des Pipeline-Patterns ist zu verstehen, dass ein
Outline-Schritt mehrere CalcJobs einreichen kann und
\texttt{ToContext} die WorkChain anschließend auf deren Abschluss
warten lässt, bevor der nächste Outline-Schritt beginnt.

\subsection{Vergleichende Beobachtungen zum Scatter-Gather-Pattern}
\label{sec:lokal:sg:zusammenfassung}

Tabelle~\ref{tab:sg:vergleich} stellt die Beobachtungen aus den fünf
Systemabschnitten entlang der Vergleichsaspekte gegenüber. Da die
beteiligten Dateien und Bausteine gegenüber dem Pipeline-Pattern
unverändert bleiben, tritt bei der Modellierung an die Stelle der
Artefaktzahl die Frage, welche Stellen eine Änderung der
Teilstückzahl betrifft. Auch die Trennung von Logik und Umgebung
sowie die verfügbaren Mittel der Nachvollziehbarkeit entsprechen dem
Pipeline-Pattern, aufgeführt ist deshalb nur, wie sich die parallele
Phase darin abbildet.

\begin{table}[H]
  \caption{Umsetzung des Scatter-Gather-Patterns im Überblick}
  \label{tab:sg:vergleich}
  \centering
  \scriptsize
  \setlength{\tabcolsep}{3pt}
  \renewcommand{\arraystretch}{1.4}
  \begin{tabular}{>{\raggedright\arraybackslash}p{2.7cm} >{\raggedright\arraybackslash}p{1.9cm} >{\raggedright\arraybackslash}p{1.9cm} >{\raggedright\arraybackslash}p{1.9cm} >{\raggedright\arraybackslash}p{1.9cm} >{\raggedright\arraybackslash}p{1.9cm}}
    \hline
     & \textbf{Nextflow} & \textbf{Stream\-Flow} & \textbf{Merlin} & \textbf{Pegasus} & \textbf{AiiDA} \\
    \hline
    Ausdruck des Patterns & \texttt{flatten} und \texttt{collect} & \texttt{scatter}-Attribut & Parameter \texttt{CHUNK}, \texttt{compute\_*} & Schleife im Erzeugungs\-skript & Schleife im Outline-Schritt, \texttt{ToContext} \\
    Änderung der Teilstückzahl & ein Parameter & Eingabewert und Namensliste in \texttt{config.yml} & zwei Stellen der Spezifikation & Aufruf\-argument, Workflow-Datei wird neu erzeugt & Aufruf\-argument, WorkChain bleibt unverändert \\
    Datenfluss und Abhängigkeiten & \texttt{path}-Angaben und Channels & Verweise auf Schritt\-ausgaben und Datei-/Feldtypen & \texttt{depends} und Pfade in den Befehlen & \texttt{File}-Objekte als Ein- und Ausgaben & Quellordner und Zuordnungen je Instanz \\
    Ausführung und Steuerung & Nextflow-Laufzeit, eine Instanz je Teilstück & StreamFlow-Laufzeit, eine Instanz je Teilstück & Celery-Worker, eine Aufgabe je Instanz & HTCondor, ein Job je Teilstück & Daemon/Runner, ein CalcJob je Teilstück \\
    Trennung Logik und Umgebung & Konfigurations\-datei unverändert & Deployment-Datei unverändert & Warteschlange unverändert & abstrakte Beschreibung und Kataloge unverändert & Rechner und Code unverändert \\
    Nachvollziehbarkeit & Instanzen als Aufgaben desselben Prozesses & Instanzen einzeln in der Terminal\-ausgabe & Instanzen einzeln in den Worker-Ausgaben & Instanzen einzeln in der Status\-abfrage & Instanzen einzeln in Status\-abfrage und Graph \\
    \hline
  \end{tabular}
\end{table}

Beim Ausdruck des Patterns unterscheiden sich die Systeme vor allem
darin, womit die parallele Phase entsteht. Nextflow und
StreamFlow drücken sie in der Workflow-Beschreibung selbst aus,
Nextflow über die Operatoren \texttt{flatten} und \texttt{collect}
auf den Channels, StreamFlow über das \texttt{scatter}-Attribut des
Verarbeitungsschritts. Merlin erreicht die Vervielfachung über die
Werteliste eines Parameters, aus der beim Übergeben des Laufs die
Verarbeitungsschritte expandiert werden. Pegasus und AiiDA verwenden
dagegen kein Sprachmittel des Systems, sondern eine Python-Schleife,
bei Pegasus im Erzeugungsskript, sodass die Verarbeitungsjobs
bereits ausgeschrieben in der abstrakten Beschreibung stehen, bei AiiDA innerhalb eines Outline-Schritts, aus dem die Instanzen
gemeinsam eingereicht werden.

Bei der Modellierung zeigt sich der deutlichste Unterschied darin,
was eine andere Teilstückzahl nach sich zieht. In Nextflow ist sie
ein Parameter, die Fassungen unterscheiden sich in einer Zeile, in
StreamFlow ändern sich der Eingabewert und die mitwachsende Liste
der Ergebnisnamen in der Konfigurationsdatei und in Merlin sind
zwei Stellen der Spezifikation aufeinander abzustimmen, die
Werteliste des Parameters und das Argument des Zerlegeschritts.
Pegasus und AiiDA behandeln die Teilstückzahl beide als
Aufrufargument, greifen damit aber an unterschiedlichen Stellen ein,
bei Pegasus im Erzeugungsskript, sodass die Workflow-Datei neu
erzeugt wird und mit der Teilstückzahl wächst, bei AiiDA im
Submit-Skript, sodass die Beschreibung unverändert bleibt.

Die Sichtbarkeit des Datenflusses behält in allen fünf Systemen die
Form des Pipeline-Patterns, die Unterschiede liegen darin, wie der
Weg der Teilstücke darin erscheint. In Nextflow ist die Ausgabe des
Zerlegens als Dateimuster deklariert, an dem erkennbar ist, dass
mehrere Dateien entstehen und jede Instanz erhält eines der
Teilstücke als deklarierte Eingabe. In StreamFlow unterscheiden die
Typangaben zwischen einer einzelnen Datei und einem Feld von
Dateien, das Zerlegen liefert ein Feld, jede Instanz erhält daraus
ein Element und das Zusammenführen erhält wieder das ganze Feld. In
Merlin stehen die Wege als Pfade im Befehl, jede Instanz liest ihr
Teilstück über den Parameterwert und das Zusammenführen sammelt die
Ergebnisse über ein Pfadmuster aus allen Arbeitsverzeichnissen ein.
In Pegasus erscheint dasselbe Dateiobjekt beim erzeugenden Job als
Ausgabe und beim lesenden als Eingabe, sodass der Weg jedes
Teilstücks über die Objekte verfolgbar ist. In AiiDA erhält jede
Instanz den Ordner des Zerlegeschritts als Quelle mit der Zuordnung
genau ihres Teilstücks, das Zusammenführen die Ordner aller
Instanzen.

Bei der Ausführung bleibt die ausführende Komponente gegenüber dem
Pipeline-Pattern unverändert, neu ist, dass sie je Teilstück eine
eigene Einheit startet, eine Prozessinstanz in Nextflow und
StreamFlow, einen expandierten Schritt in Merlin, einen Job in
Pegasus und einen CalcJob in AiiDA. In allen fünf Systemen liefen
in den durchgeführten Läufen diese Einheiten nebeneinander und das Zusammenführen begann erst, wenn alle abgeschlossen waren.

Die Trennung von Workflow-Logik und Ausführungsumgebung bleibt in
allen fünf Systemen unverändert bestehen, die parallele Phase
berührt sie nicht. Der Mechanismus, der die Trennung trägt, ist
derselbe wie beim Pipeline-Pattern, die Konfigurationsdatei in
Nextflow, die Deployment-Datei in StreamFlow, die Warteschlange in
Merlin, die abstrakte Beschreibung mit den Katalogen in Pegasus und
die als Eingaben behandelten Objekte für Rechner und Code in AiiDA.

Bei der Nachvollziehbarkeit stellt sich mit der parallelen Phase
eine Frage, die das Pipeline-Pattern nicht kannte, nämlich ob die
Verarbeitungsinstanzen einzeln sichtbar werden. Am einen Ende steht
Nextflow, dessen Fortschrittsanzeige die Instanzen als Aufgaben
desselben Prozesses zählt, ohne sie einzeln zu benennen. In
StreamFlow, Merlin und Pegasus erscheinen die Instanzen einzeln,
in der Terminalausgabe, in den Ausgaben des Workers beziehungsweise
in der Statusabfrage des Laufs. In AiiDA sind sie über die Statusabfrage hinaus als eigene Prozessknoten im Provenance-Graphen der Datenbank festgehalten,
sodass die Zuordnung der Instanzen auch nach dem Lauf abfragbar
bleibt.

\section{Ergebnisse der lokalen Messungen}
\label{sec:lokal:messungen}

Die folgenden Ergebnisse beruhen auf fünf Wiederholungen je
Konfiguration bei der kurzen und der mittleren Rechenlast sowie drei
Wiederholungen bei der langen, berichtet wird jeweils der Median.
Für jede Konfiguration wurde zusätzlich ein Referenzlauf erhoben,
der dieselben Benchmark-Skripte in derselben Reihenfolge ohne
Workflow-Management-System ausführt, jeder Wiederholung ist dabei
ein eigener Referenzlauf zugeordnet. Als Overhead wird im Folgenden
der Median der je Wiederholung gebildeten Differenzen zwischen der
Gesamtlaufzeit des Systemlaufs und der des zugehörigen
Referenzlaufs bezeichnet. Er beschreibt den zusätzlich
beobachteten Zeitbedarf gegenüber der Ausführung derselben
Rechenlogik ohne Workflow-Management-System. Neben den Gesamtlaufzeiten
werden zwei Koordinationsgrößen betrachtet, die Übergangslatenz
zwischen dem Ende eines Schritts und dem Beginn des nächsten sowie
die Startspreizung der parallelen Verarbeitungsinstanzen, also der
Abstand zwischen dem ersten und dem letzten Instanzstart.

Abbildung~\ref{fig:mess:pipe} zeigt die Gesamtlaufzeiten des
Pipeline-Patterns je Rechenlast. Vier Systeme liegen dabei in einem
engen Band, bei der kurzen Rechenlast zwischen 5,8 und
9,2\,Sekunden, bei der mittleren zwischen 53,4 und 63,1\,Sekunden
und bei der langen zwischen 238 und 280\,Sekunden, wobei die
Reihenfolge innerhalb des Bandes zwischen den Rechenlasten wechselt.
Pegasus liegt mit 184, 228 und 442\,Sekunden durchgängig deutlich
darüber. Sein Abstand zum eigenen Referenzlauf beträgt bei allen
drei Rechenlasten zwischen 171 und 180\,Sekunden und ist damit
nahezu konstant, was dem vorgelagerten Planungsschritt und der
Ausführung über HTCondor entspricht, die unabhängig von der
Rechenlast anfallen. Auffällig ist zudem AiiDA bei der kurzen
Rechenlast, dessen Laufzeit von 9,2\,Sekunden erkennbar über den
drei übrigen Systemen liegt, während es bei der mittleren und der
langen Rechenlast im Band bleibt.

\begin{figure}[H]
  \centering
  \includegraphics[width=\textwidth]{content/figs/lokal_mac_graphen/02_makespan_pipeline}
  \caption{Median der Gesamtlaufzeit des Pipeline-Patterns je Rechenlast auf der lokalen Ebene, logarithmische Skala}
  \label{fig:mess:pipe}
\end{figure}

Abbildung~\ref{fig:mess:sg} zeigt die Gesamtlaufzeiten des
Scatter-Gather-Patterns über die Zahl der Teilstücke. Bei der
mittleren und der langen Rechenlast sinken die Laufzeiten der vier
Systeme im engen Band mit jeder Verdopplung der Teilstückzahl
deutlich, bei der langen Rechenlast etwa von 248 auf 66\,Sekunden in
StreamFlow und von 255 auf 73\,Sekunden in AiiDA. Bei der kurzen
Rechenlast fällt der Rückgang geringer aus, AiiDA liegt dort bei
vier Teilstücken mit 10,7\,Sekunden sogar leicht über dem Wert bei
zwei Teilstücken. Pegasus folgt dem Verlauf ebenfalls, jedoch
flacher, seine Laufzeiten sinken bei der langen Rechenlast von 480
auf 289\,Sekunden, bei der kurzen bleiben sie mit rund
207\,Sekunden über alle Teilstückzahlen nahezu unverändert. Der
Anteil, der nicht auf die Rechenlogik zurückgeht, verändert sich
dabei in allen Systemen nur wenig, in StreamFlow, Nextflow und
Merlin liegt er bei der langen Rechenlast zwischen 1,5 und
3,5\,Sekunden, in AiiDA zwischen 8,9 und 10,3\,Sekunden und in
Pegasus zwischen 216 und 221\,Sekunden.

\begin{figure}[H]
  \centering
  \includegraphics[width=\textwidth]{content/figs/lokal_mac_graphen/03_makespan_scatter_gather}
  \caption{Median der Gesamtlaufzeit des Scatter-Gather-Patterns über die Zahl der Teilstücke je Rechenlast auf der lokalen Ebene, logarithmische Skala}
  \label{fig:mess:sg}
\end{figure}

Abbildung~\ref{fig:mess:overhead} stellt den Overhead aller zwölf
Konfigurationen gegenüber. Die Systeme bilden drei Gruppen, deren
Reihenfolge in jeder einzelnen Konfiguration gleich bleibt.
StreamFlow, Nextflow und Merlin liegen zwischen 1,4 und
4,2\,Sekunden. AiiDA liegt zwischen 5,1 und 10,3\,Sekunden und
Pegasus zwischen 171 und 221\,Sekunden. Ein proportionaler
Zusammenhang zwischen Rechenlast und Overhead ist in keiner der
Gruppen erkennbar. Ein Zusammenhang mit dem Umfang des Workflows
zeigt sich dagegen bei einzelnen Systemen. Beim
Scatter-Gather-Pattern mit seinen sechs Schritten liegt der
Overhead von Pegasus durchgängig und der von AiiDA bei der kurzen
und der mittleren Rechenlast über dem des Pipeline-Patterns mit
vier Schritten. Bei AiiDA steht dies im Einklang mit der
Abwicklung jedes Schritts als eigener CalcJob über Daemon und
Datenbank, bei Pegasus mit den zusätzlichen Übergängen des Musters, an denen jeweils der Abfragezyklus anfällt.

\begin{figure}[H]
  \centering
  \includegraphics[width=\textwidth]{content/figs/lokal_mac_graphen/01_overhead_alle}
  \caption{Overhead gegenüber dem Referenzlauf für alle zwölf Konfigurationen auf der lokalen Ebene, logarithmische Skala}
  \label{fig:mess:overhead}
\end{figure}

Abbildung~\ref{fig:mess:latenz} zeigt die Übergangslatenzen zwischen
aufeinanderfolgenden Schritten. Die Systeme unterscheiden sich darin
über mehrere Größenordnungen, in StreamFlow liegt der Median je
Übergang bei etwa 0,06\,Sekunden, in Nextflow bei 0,09, in Merlin
bei 0,38, in AiiDA bei 1,16 und in Pegasus bei 20,1\,Sekunden.
Innerhalb eines Systems sind die Werte über die einzelnen Übergänge
und über beide Muster hinweg weitgehend gleich. Die Übergangslatenz zeigt sich damit im lokalen Aufbau als eine weitgehend systemspezifische Größe. Sie wird maßgeblich von der Stelle bestimmt, an der ein Schritt an den nächsten übergeben wird, in Merlin von der Warteschlange, in AiiDA vom Daemon mit der Datenbank und in Pegasus vom Abfragezyklus, der die Bereitschaft der Jobs prüft. In Nextflow und StreamFlow startet die Laufzeit die Schritte selbst, entsprechend klein bleibt der Abstand.

\begin{figure}[H]
  \centering
  \includegraphics[width=\textwidth]{content/figs/lokal_mac_graphen/05_lokal_uebergangslatenzen}
  \caption{Median der Übergangslatenzen zwischen aufeinanderfolgenden Schritten auf der lokalen Ebene, logarithmische Skala}
  \label{fig:mess:latenz}
\end{figure}

Abbildung~\ref{fig:mess:spreizung} zeigt die Startspreizung der
parallelen Verarbeitungsinstanzen, also den Abstand zwischen dem
ersten und dem letzten Instanzstart. Die Werte reichen dabei von wenigen Millisekunden bis in den Bereich einer knappen Sekunde. Bei zwei Teilstücken reicht
der Median von 1\,Millisekunde in Nextflow über 9 in StreamFlow und
19 in Pegasus bis zu 58 in AiiDA und 228 in Merlin, bei vier
Teilstücken von 7 über 40 und 62 bis zu 698 und 734\,Millisekunden.
In allen fünf Systemen wächst die Spreizung mit der Zahl der
Teilstücke, am deutlichsten in AiiDA, wo sie sich von zwei auf vier
Teilstücke etwa verzwölffacht. Die Staffelung entspricht dem Weg, über den die Instanzen untereinander gestartet werden. In Nextflow, StreamFlow und Pegasus
entstehen sie ohne Zwischenstelle, in Nextflow und StreamFlow startet
die Laufzeit sie unmittelbar nacheinander, in Pegasus liegen die
Jobs bereits geplant vor und werden gemeinsam in die Warteschlange
gegeben. In AiiDA und Merlin wird dagegen jede Instanz einzeln in eine
Warteschlange eingereiht, aus der sie in AiiDA ein Daemon-Runner
und in Merlin ein Worker entgegennimmt, und der Abstand wächst dort
mit jeder weiteren Instanz.

\begin{figure}[!htbp]
  \centering
  \includegraphics[width=\textwidth]{content/figs/lokal_mac_graphen/04_startspreizung}
  \caption{Median der Startspreizung der parallelen Verarbeitungsinstanzen auf der lokalen Ebene, logarithmische Skala}
  \label{fig:mess:spreizung}
\end{figure}

Über die einzelnen Kenngrößen hinweg fügen sich die Befunde zu
einem stimmigen Bild. Bei Pegasus entspricht der höhere Overhead
des Scatter-Gather-Patterns näherungsweise den zusätzlichen
Übergangslatenzen, die Differenz der Overhead-Mediane beider Muster beträgt
37,9\,Sekunden, während das Muster zwei Übergänge mehr enthält als
die Pipeline und zweimal 20,1\,Sekunden gemessener Latenz
40,2\,Sekunden ergeben. Zugleich zeigt Pegasus, dass die beiden
Koordinationsgrößen verschiedene Vorgänge erfassen, seine Übergänge
sind mit Abstand die teuersten, seine Startspreizung liegt mit 19
und 62\,Millisekunden dagegen im unteren Feld. AiiDA und Merlin
wiederum, die als einzige jede Instanz einzeln in eine
Warteschlange einreihen, liegen in beiden Koordinationsgrößen über
Nextflow und StreamFlow, deren Laufzeitumgebung die Schritte und
Instanzen unmittelbar selbst startet. Ein unmittelbarer Schluss von
den Koordinationsgrößen auf die Höhe des Overheads ist damit nicht
möglich. Merlin weist höhere Übergangslatenzen auf als Nextflow und
StreamFlow, liegt beim Overhead aber in derselben Gruppe und in
einzelnen Konfigurationen darunter. Die Koordinationsgrößen zeigen,
an welchen Stellen ein System zusätzliche Zeit beansprucht, sie
bestimmen den Overhead jedoch nicht allein.

\section{Zwischenfazit}
\label{sec:lokal:zwischenfazit}

Beide Workflow-Muster ließen sich in allen fünf Systemen mit
denselben Benchmark-Skripten umsetzen, die Unterschiede liegen in
der Beschreibung, der Ausführung und der Art, wie ein Lauf
nachvollziehbar wird. Bei der Modellierung trennen sich die Systeme
darin, ob die Reihenfolge aus den Datenbeziehungen entsteht oder
ausdrücklich benannt wird, ob eine eigene Laufzeit ausführt oder
eine eigenständige Komponente, und wann die auszuführenden Aufgaben
feststehen. Eine geänderte Rechenlast betrifft in allen Systemen
nur eine Stelle, eine geänderte Teilstückzahl reicht dagegen von
einem einzelnen Parameter bis zum erneuten Erzeugen der
Workflow-Beschreibung.

Die Messungen zeigen drei Overhead-Gruppen, deren Reihenfolge in
allen zwölf Konfigurationen gleich bleibt, StreamFlow, Nextflow und
Merlin im unteren Bereich, AiiDA darüber und Pegasus mit deutlichem
Abstand darüber. Ein proportionaler Zusammenhang zwischen
Rechenlast und Overhead ist dabei nicht erkennbar. Unterschiede
zwischen den beiden Workflow-Mustern zeigen sich insbesondere bei
AiiDA und Pegasus. Die Übergangslatenz reicht über mehrere
Größenordnungen von StreamFlow bis Pegasus und bleibt innerhalb
eines Systems weitgehend stabil, während die Startspreizung mit der
Zahl der parallelen Instanzen wächst. Die Koordinationsgrößen
machen damit Unterschiede in der Übergabe aufeinanderfolgender
Schritte und im Start paralleler Instanzen sichtbar, ohne die Höhe
des gesamten Overheads allein zu bestimmen.

Diese Befunde beschreiben das Verhalten der Systeme unter den
kontrollierten Bedingungen eines einzelnen Rechners.
Kapitel~\ref{sec:mogon} untersucht, wie sich die Systeme auf einem
Slurm-basierten HPC-System bereitstellen und betreiben lassen.